% Chapter 13: maintenance.
\chapter{Maintenance}\label{ch:maintenance}
\index{maintenance}

The controller needs no routine service. The parts that matter for accurate
rolls are the measuring wheel and the paper path. Always switch the power
off first.

\begin{center}
\renewcommand{\arraystretch}{1.45}
\begin{tabularx}{\textwidth}{@{}>{\raggedright\arraybackslash}p{0.20\textwidth}>{\raggedright\arraybackslash}X@{}}
\toprule
\textbf{How often} & \textbf{Task} \\
\midrule
Every shift & Look at the measuring wheel: it must sit on the paper and turn
  freely. Wipe off paper dust with a dry cloth. Glance at the \shows{DONE}
  overrun on the first roll. \\
Weekly & Clean the wheel's rubber with a lightly damp cloth; let it dry. Check
  that the wheel's spring or weight still presses it onto the paper. Check the
  encoder cable is plugged in and not rubbing. \\
Monthly & Measure one 30 yard roll by hand and write it in the log below. Check
  that all four buttons click and all four LEDs light (idle shows \led{1} and
  \led{3}; manual shows \led{4}). \\
When rolls drift & Recalibrate (Chapter~\ref{ch:calibration}). \\
When the rubber is worn, cracked or glazed & Replace the wheel
  ([SPARE PART NUMBER]), then calibrate with the length check. \\
\bottomrule
\end{tabularx}
\end{center}

\begin{note}
Do not oil the measuring wheel or the paper path. Oil makes the wheel slip,
and a slipping wheel winds long or short rolls without any warning.
\end{note}

\section*{Maintenance log}
\begin{center}
\renewcommand{\arraystretch}{2.0}
\begin{tabularx}{\textwidth}{|p{0.15\textwidth}|p{0.15\textwidth}|X|p{0.12\textwidth}|}
\hline
\textbf{Date} & \textbf{Roll measured} & \textbf{Work done / notes} & \textbf{Initials} \\
\hline
 & & & \\ \hline
 & & & \\ \hline
 & & & \\ \hline
 & & & \\ \hline
 & & & \\ \hline
 & & & \\ \hline
 & & & \\ \hline
 & & & \\ \hline
\end{tabularx}
\end{center}
